\chapter{Referência Rápida}
\label{apendice:referencia}

\section{Tipos}

\begin{center}
\begin{tabular}{lll}
\toprule
\textbf{Tipo} & \textbf{Exemplo} & \textbf{Notas} \\
\midrule
\texttt{Int}       & \texttt{42}, \texttt{-7}      & 64 bits, com sinal \\
\texttt{Float}     & \texttt{3.14}, \texttt{1e-5}  & IEEE 754 \\
\texttt{Bool}      & \texttt{true}, \texttt{false} & \\
\texttt{Nil}       & \texttt{nil}                  & Ausência de valor \\
\texttt{String}    & \texttt{"texto"}              & UTF-8 \\
\texttt{FString}   & \texttt{f"val: \{x\}"}        & Interpolação \\
\texttt{List}      & \texttt{[1, 2, 3]}            & Índice base 0 \\
\texttt{Dict}      & \texttt{\{"a": 1\}}           & Chave string ou int \\
\texttt{DataFrame} & \texttt{load "f.csv" as df}   & COW, colunas Float \\
\bottomrule
\end{tabular}
\end{center}

\section{Operadores}

\begin{center}
\begin{tabular}{ll}
\toprule
\textbf{Aritmético}    & \texttt{+  -  *  /  \^{}  \%} \\
\textbf{Comparação}    & \texttt{== != > < >= <=} \\
\textbf{Lógico}        & \texttt{and  or  not} \\
\textbf{Pertencimento} & \texttt{in} \\
\textbf{Pipe}          & \texttt{|>} \\
\textbf{Atribuição}    & \texttt{=  +=  -=  *=  /=} \\
\textbf{Range}         & \texttt{a..b} (exclusivo) \quad \texttt{a..=b} (inclusivo) \\
\bottomrule
\end{tabular}
\end{center}

\section{Ranges}

Ranges são expressões que produzem listas de inteiros:

\begin{center}
\begin{tabular}{ll}
\toprule
\texttt{1..5}          & \texttt{[1, 2, 3, 4]} \\
\texttt{1..=5}         & \texttt{[1, 2, 3, 4, 5]} \\
\texttt{5..1}          & \texttt{[5, 4, 3, 2]} (decrescente) \\
\texttt{chain(s1, s2)} & Concatena listas/ranges \\
\bottomrule
\end{tabular}
\end{center}

\begin{lstlisting}[style=inline]
chain(1..3, 9..=11)           // [1, 2, 9, 10, 11]
map(1..=5, fn(x) { x^2 })    // [1, 4, 9, 16, 25]
\end{lstlisting}

\section{Palavras-chave}

\begin{lstlisting}[style=inline]
let  const  fn  return  if  else  for  in  while  break  continue
match  try  catch  and  or  not  true  false  nil
\end{lstlisting}

\section{Estatística descritiva e momentos}

Todas as funções abaixo aceitam lista pura ou coluna de DataFrame.
As formas sobre DataFrame aceitam \lstinline{if = cond}.

\begin{center}
\begin{tabular}{lll}
\toprule
\textbf{Notação} & \textbf{Função} & \textbf{Formas} \\
\midrule
$E(X)$
  & \texttt{mean(df, x)}
  & lista / df / \texttt{if =} \\
$\mathrm{Var}(X)$
  & \texttt{variance(df, x)}
  & lista / df / \texttt{if =} \\
$\mathrm{DP}(X)$
  & \texttt{sd(df, x)}
  & lista / df / \texttt{if =} \\
$\mathrm{Cov}(X,Y)$
  & \texttt{cov(df, x, y)}
  & df / \texttt{if =} \\
$\rho(X,Y)$
  & \texttt{corr\_pair(df, x, y)}
  & df / \texttt{if =} \\
$\mathrm{Med}(X)$
  & \texttt{median(df, x)}
  & lista / df / \texttt{if =} \\
$Q_p(X)$
  & \texttt{quantile(df, x, p)}
  & lista / df / \texttt{if =} \\
\midrule
\multicolumn{2}{l}{\texttt{correlate(df, x, y, ...)}}
  & matriz de correlação (display) \\
\multicolumn{2}{l}{\texttt{summarize(df, x)}}
  & dict: mean, sd, variance, min, max, N \\
\multicolumn{2}{l}{\texttt{summarize(df, x, detail=true)}}
  & + p1..p99, skew, kurt \\
\bottomrule
\end{tabular}
\end{center}

\begin{lstlisting}[style=inline]
mean(df, renda, if = sexo == 1)         // E(renda | homens)
cov(df, educ, salario, if = formal==1)  // Cov condicional
quantile(df, renda, 0.9)               // percentil 90
\end{lstlisting}

\section{Funções de dados}

\begin{center}
\begin{tabular}{ll}
\toprule
\textbf{Função} & \textbf{Descrição} \\
\midrule
\texttt{load "f" as df}           & Carrega arquivo \\
\texttt{save df "f"}              & Salva arquivo \\
\texttt{describe df}              & Estrutura do DataFrame \\
\texttt{count(df)}                & Número de linhas \\
\texttt{count(df, cond)}          & Linhas onde cond é verdade \\
\texttt{nrow(df)}                 & Alias de count(df) \\
\texttt{filter(df, cond)}         & Filtra linhas \\
\texttt{select(df, cols...)}      & Seleciona colunas \\
\texttt{drop(df, cols...)}        & Remove colunas \\
\texttt{mutate(df, col=expr)}     & Adiciona/modifica coluna \\
\texttt{generate(df, col=expr)}   & Alias de mutate \\
\texttt{sort(df, col [desc])}     & Ordena \\
\texttt{dropna(df [, cols...])}   & Remove NaN \\
\texttt{append(df1, df2)}         & Empilha verticalmente \\
\texttt{group\_by(df, col, ...)}  & Agrega por grupo \\
\texttt{pivot\_longer(...)}       & Largo → longo \\
\texttt{pivot\_wider(...)}        & Longo → largo \\
\texttt{summarize(df, col)}       & Estatísticas descritivas \\
\texttt{codebook df}              & Sumário completo \\
\texttt{tabulate df col}          & Frequências \\
\bottomrule
\end{tabular}
\end{center}

\section{Operadores de séries temporais}

Disponíveis dentro de \lstinline{generate df col = expr}:

\begin{center}
\begin{tabular}{ll}
\toprule
\texttt{L.x}   & Lag de ordem 1 \\
\texttt{L2.x}  & Lag de ordem 2 \\
\texttt{F.x}   & Lead de ordem 1 \\
\texttt{D.x}   & Diferença primeira \\
\bottomrule
\end{tabular}
\end{center}

\section{Atalhos do REPL}

\begin{center}
\begin{tabular}{ll}
\toprule
\textbf{Atalho}  & \textbf{Ação} \\
\midrule
\texttt{exit}    & Sair do REPL \\
\texttt{Ctrl-D}  & Sair do REPL (EOF) \\
\texttt{Ctrl-C}  & Cancelar linha atual \\
\texttt{Ctrl-R}  & Buscar no histórico \\
\bottomrule
\end{tabular}
\end{center}
